% ============================================================
\section{Background and Related Work}
\label{sec:related}
% ============================================================

\subsection{Graph-Based Extractive Summarization}
Extractive summarization preserves source wording and requires no summary-specific
training data. Graph-based methods represent sentences as nodes and their relations
as weighted edges, then estimate sentence importance from the graph structure.
TextRank applies PageRank to this representation~\cite{mihalcea2004textrank}, but
its performance depends on the edge-weight function. Word overlap and variants
based on TF-IDF, LexRank, BM25+, or LDA-extended TextRank primarily capture lexical
relations~\cite{vora2024extractive}. Sentence embeddings instead model semantic
similarity between lexically different sentences~\cite{kadriu2021extractive}.

\subsection{Keyphrase Extraction}
Unsupervised keyphrase extraction includes statistical methods such as RAKE
~\cite{rose2010automatic} and YAKE~\cite{campos2020yake}; graph-based methods such
as TextRank~\cite{mihalcea2004textrank}, TopicRank~\cite{bougouin2013topicrank},
MultipartiteRank~\cite{boudin2018unsupervised}, and WordAttractionRank
~\cite{wang2014corpus}; and embedding-based methods. KeyBERT
~\cite{grootendorst2020keybert} ranks candidate n-grams by cosine similarity to the
document and can apply maximal marginal relevance (MMR)~\cite{carbonell1998use} to
reduce redundancy. Here, matching a short phrase to a document summary is an
asymmetric retrieval task; a retrieval-oriented embedding model is therefore more
appropriate than one optimized only for symmetric semantic textual similarity
(STS).

\subsection{Label-Space Discovery}
Label-space discovery has used topic models, embedding-and-clustering methods, and
LLM-assisted topic naming. Many such pipelines assume a fixed topic count or a
single dominant topic per document. X-MLClass~\cite{li2024open} uses an LLM for
chunking, keyphrase generation, clustering, and iterative rare-label discovery.
However, repeated raw-text processing increases token usage and latency and may
produce entity-level or journalistic labels on news data. StatLM instead places a
statistical encoder before the LLM so that label discovery operates on condensed
keyphrases.

\subsection{Vietnamese Processing Resources}
StatLM uses pyvi~\cite{pyvi} for word segmentation and VnCoreNLP
~\cite{vu2018vncorenlp} for segmentation, part-of-speech (POS) tagging, and
named-entity recognition (NER). The encoder models are matched to their tasks
(Table~\ref{tab:emb}). \emph{dangvantuan/vietnamese-embedding}
~\cite{dangvantuan2021vietnamese}, a Sentence-BERT model
~\cite{reimers2019sentence} based on PhoBERT~\cite{nguyen2020phobert}, is optimized
for symmetric STS and reports Pearson 84.21 on \emph{vi-stsbenchmark}.
\emph{bkai-foundation-models/vietnamese-bi-encoder}
~\cite{nguyen2025improving} is trained for asymmetric retrieval and is reported to
rank first on the Vietnamese VCS benchmark.
